\begin{tikzpicture}[
  font=\scriptsize,
  execute at begin node={\hyphenpenalty=10000\exhyphenpenalty=10000},
  aktor/.style={draw, rounded corners, fill=blue!8, minimum width=2.6cm, minimum height=0.8cm, align=center},
  garis/.style={dashed, gray},
  ket/.style={font=\tiny, align=center},
  catat/.style={draw, rounded corners, fill=orange!12, font=\tiny, align=center, text width=2.5cm},
  panah/.style={-{Stealth[length=2mm]}, thick},
  balik/.style={-{Stealth[length=2mm]}, thick, dashed}
]
  \node[aktor] (lab) at (0,0) {Sistem laboratorium\\{\tiny\mdseries klien mitra}};
  \node[aktor] (sa) at (4.7,0) {Server otorisasi};
  \node[aktor] (api) at (9.4,0) {API hasil laboratorium\\{\tiny\mdseries server sumber daya}};
  \draw[garis] (lab.south) -- (0,-4.4);
  \draw[garis] (sa.south) -- (4.7,-4.4);
  \draw[garis] (api.south) -- (9.4,-4.4);

  \draw[panah] (0,-1.05) -- node[above, ket] {1. identitas klien dan \textit{scope} yang diminta} (4.7,-1.05);
  \draw[balik] (4.7,-1.65) -- node[above, ket] {2. token akses berlaku 15 menit} (0,-1.65);
  \draw[panah] (0,-2.3) -- node[above, ket, pos=0.25] {3. kirim hasil dengan token} (9.4,-2.3);
  \node[catat] at (8.0,-3.05) {4. periksa token, \textit{audience}, \textit{scope}, dan hak atas pasien rujukan};
  \draw[balik] (9.4,-3.85) -- node[above, ket, pos=0.75] {5. diterima, atau ditolak dengan 403} (0,-3.85);

  \node[ket] at (4.7,-4.7) {Tidak ada pasien yang masuk dan tidak ada layar persetujuan. Klien bertindak untuk kepentingannya sendiri};
\end{tikzpicture}
